% Sub-stage F1: the tables and the moduli: the table of a prime as examine_prime builds it, with
% the Barrett reduction on the fixed schedule of Horner's rule, handing back its coefficients modulo
% p and its entry at infinity; the loop of the tables of sieving_info, with the exit and the
% forbidden divisors prime by prime; prime powers and products of them as moduli, with the sieve
% entry of a modulus of any kind; and the parameters that admit them.
\chapter{The tables and the moduli}\label{chap:tables}

Stage~F models the side of the numerators: the tables of the moduli, the three stages of the sieve,
the packing of the numerators of a class by its stride, and the first stage in chunks of bit arrays
with the rows as the builders of ratpoints lay them out. It runs in four sub-stages with a chapter
each: F1, the tables and the moduli (this chapter); F2, the stages of the sieve and the check; F3,
the packing; F4, the chunked first stage and the faithful rows. Each of them only computes what the
model computes already in another way, or removes candidates that are not points, so the points
found and their order do not change.

F1 models what ratpoints does with its tables before the sieve. The table of a prime is computed as
\verb|examine_prime| computes it, by Horner's rule in the arithmetic of 64-bit words with the
Barrett reduction on a fixed schedule, and it hands back to its caller the coefficients modulo the
prime and its entry at infinity; the forbidden divisors of Chapter~\ref{chap:denominators} take them
from it, prime by prime in the loop over the primes that \verb|sieving_info| examines, which ends
the search at a prime modulo which the curve has no points. And the moduli may be prime powers and
products of prime powers, as those of ratpoints are: their tables are those of \verb|examine_power|,
and their rows are the rows of their factors, combined.

In this sub-stage the theorems of Chapter~\ref{chap:denominators}, the refusal theorems and the
statements that comparator certifies keep their text, and change only through the parameters that
they quantify over (Definition~\ref{def:Params}): the parameters have the number of primes that the
loop of the tables examines, and their well-formedness admits composite moduli, below
$2^{\mathrm{maxBitsInPrime}}$ as in the C code, where before it admitted the $171$ primes of the
table; it does not admit the primes of the table from $2^{\mathrm{maxBitsInPrime}}$ on. The bound is
the C code's own: ratpoints examines only the primes below $2^{\mathrm{maxBitsInPrime}}$, and its
schedule of the Barrett reduction is exact for them (and, at the $8$ bits of the build, for no
larger prime from degree $7$ on). The parameters of the command line
(Definition~\ref{def:sieveParams}) are well-formed as before, and the output is the same for every
well-formed choice.

\section{What the C code does with its tables}\label{sec:tablesC}

For the polynomial searched, of degree $d$ with the coefficients $c_0, \dots, c_d$ and the form
$F(x, z) = \sum_k c_k x^k z^{D - k}$ of even degree $D = d + (d \bmod 2)$, \verb|sieving_info|:
\begin{enumerate}
  \item examines the first \verb|pn_lim| primes of the table in turn, \verb|pn_lim| being at first
    the number of primes asked for, \verb|num_primes|, at most $\mathrm{numPrimes}$, the primes
    below $2^{\mathrm{maxBitsInPrime}}$, and raised when the number was left to the library and the
    choice of the moduli needs more primes (\csource{find_points.c}{2605}{2707}).
    \verb|examine_prime| (\csource{find_points.c}{1683}{1826}) reduces the coefficients modulo $p$,
    by \verb|mod| for \verb|long|s and by the Euclidean remainder for GMP integers, into an array of
    its caller; sets the entry $0$ of the table from the squares modulo $p$ at $c_0$, and the entry
    $a$ for $0 < a < p$ from $f(a)$, computed by Horner's rule in \verb|unsigned long|s with the
    Barrett reduction, once at the end for $d$ at most
    $\mathrm{RP\_HORNER\_STEPS} = 64 / \mathrm{maxBitsInPrime} - 1$, and also after every
    $\mathrm{RP\_HORNER\_STEPS}$ steps otherwise (\csource{find_points.c}{1628}{1655},
    \csource{find_points.c}{1716}{1747}); sets the entry $p$, whether there is a point at infinity
    modulo $p$; and returns $-1$, the exit, when no entry is set, $0$ when every residue gives a
    square and there is no point at infinity, and otherwise makes the sieve entry of the prime. It
    hands back the table and the coefficients;
  \item for each examined prime, in the same loop and whatever \verb|examine_prime| returned short
    of the exit, at which \verb|sieving_info| returns at once, while the checked loop over the
    denominators is on and fewer than \verb|max_forbidden| forbidden divisors have been found,
    records the forbidden divisor of $p$ from the coefficients modulo $p$ and the entry at infinity
    that it handed back (\csource{find_points.c}{2624}{2677}); after the loop, the further primes of
    the table from the final position \verb|pn_lim| on (\csource{find_points.c}{2709}{2789});
  \item offers as moduli the odd composites $m$ with $9 \le m < 2^{\mathrm{maxBitsInPrime}}$ and
    $m \le \mathrm{RATPOINTS\_COMPOSITE\_MAX}$, $64$ by default, that factor over the examined
    primes into at most three prime powers, each of which says something about the curve, and whose
    table is not too dear for the run, while there is room for its prime powers and for more
    candidates (\verb|add_moduli|, \csource{find_points.c}{2339}{2411}); for a factor $p^e$ with
    $e \ge 2$, \verb|examine_power| (\csource{find_points.c}{2254}{2331}) computes the coefficients
    modulo $q = p^e$, the squares modulo $q$, for each residue $x$ whether $f(x)$ is a square modulo
    $q$, for each multiple $t$ of $p$ whether the reversed polynomial
    $\tilde f(t) = \sum_k c_k t^{D - k}$ is one, and the inverses modulo $q$ of the units, by the
    extended Euclidean algorithm (\verb|modinv|, \csource{find_points.c}{2241}{2252});
  \item ranks the primes and the composites by what they say per cost and takes the moduli of the
    first stage in that order, then ranks the rest again with a key for the second stage and takes
    its moduli, pairwise coprime throughout, and makes the sieve entry of each composite
    (\verb|make_power| and \verb|make_modulus|, \csource{find_points.c}{2414}{2475}); the third
    stage takes primes only, and examines further primes of the table when it needs them and the
    number of primes was left to the library, with the same exit
    (\csource{find_points.c}{2976}{3000}).
\end{enumerate}
The row of a denominator $b$ in the table of a modulus is then, for each kind, a pattern of the
numerators modulo the modulus: for a prime $p$ and $r = b \bmod p \ne 0$, the residues $x$ with
$f(x r^{-1})$ a square modulo $p$, and for $r = 0$ those not divisible by $p$; for a prime power
$q = p^e$ and $r = b \bmod q$, the residues $x$ with $f(x r^{-1})$ a square modulo $q$ when
$p \nmid r$, and with $p \nmid x$ and $\tilde f(r x^{-1})$ a square modulo $q$ when $p \mid r$
(\verb|_ratpoints_sieve_init_power|, \csource{init.c}{359}{393}); for a product, the residues that
the row of every factor admits at its own residue (\verb|_ratpoints_sieve_init_product|,
\csource{init.c}{399}{445}). Every point $(a, y, b)$ with $\gcd(a, b) = 1$ passes every row: for
$p \nmid b$, $F(a, b) \equiv b^D f(a b^{-1})$ modulo $q$ with $b^D$ the square of a unit; for
$p \mid b$, $p \nmid a$ and $F(a, b) \equiv a^D \tilde f(b a^{-1})$ modulo $q$. Which moduli are
chosen, of which kinds and in which stages, changes neither the points nor their order.

\section{The table of a prime as built}\label{sec:tableBuilt}

\begin{definition}[The Barrett reduction on the schedule of Horner's rule]\label{def:barrett}
  \uses{def:primeTable}
  \lean{Ratpoints.maxPrimeEven, Ratpoints.hornerSteps, Ratpoints.one_le_hornerSteps, Ratpoints.maxBitsInPrime_le_ten, Ratpoints.maxPrimeEven_le_two_pow_10, Ratpoints.maxPrimeEven_le_two_pow_32, Ratpoints.getElem!_primeTable_lt_maxPrimeEven, List.SortedLE.getElem_lt_of_lt_length_filter, Ratpoints.Sieve.barrett, Ratpoints.Sieve.hornerBarrett}
  \leanok
  The power of two above the moduli, $\mathrm{maxPrimeEven} = 2^{\mathrm{maxBitsInPrime}}$
  (\csource{rp-private.h}{325}{361}), and the number of Horner steps between two reductions,
  $\mathrm{hornerSteps} = 64 / \mathrm{maxBitsInPrime} - 1$ (\csource{find_points.c}{1652}{1655}),
  both computed from the constant $\mathrm{maxBitsInPrime}$; two facts about the constant, checked
  by evaluation as the C code fixes them at compile time: $1 \le \mathrm{hornerSteps}$, which the
  schedule needs, and $\mathrm{maxBitsInPrime} \le 10$, since the header knows only the sizes from
  $5$ to $10$ bits and takes $7$ for any other, so that the moduli are below $2^{10}$, as the
  comments of \verb|examine_power| say, and below $2^{32}$; and the first $\mathrm{numPrimes}$
  primes of the table are below $\mathrm{maxPrimeEven}$, since the table is sorted. The Barrett
  reduction of $u$ modulo $p$ with the reciprocal $m = \lfloor (2^{64} - 1) / p \rfloor$:
  $r = u - \lfloor u m / 2^{64} \rfloor p$, and $r - p$ if $r \ge p$
  (\csource{find_points.c}{1639}{1650}, the form with a 128-bit product, which the default build and
  the build of the model compile). The value of $f$ at $a$ modulo $p$ from the coefficients modulo
  $p$ by Horner's rule from $c_d$, without a reduction for $d \le \mathrm{hornerSteps}$, and with a
  reduction after every $\mathrm{hornerSteps}$ steps otherwise, and one at the end
  (\csource{find_points.c}{1724}{1724}, \csource{find_points.c}{1726}{1731},
  \csource{find_points.c}{1735}{1735}, \csource{find_points.c}{1737}{1744}).
\end{definition}

\begin{lemma}[The Barrett reduction and the schedule of Horner's rule]\label{lem:barrett}
  \uses{def:barrett, def:examinePrime}
  \lean{Ratpoints.Sieve.barrett_eq, Ratpoints.hornerSteps_add_one, Ratpoints.pow_hornerSteps_add_one_lt, Ratpoints.Sieve.HornerInvariant, Ratpoints.Sieve.HornerInvariant.lt_two_pow_64, Ratpoints.Sieve.hornerBarrett_spec, Nat.mul_sub_one_div_div_mul_le, Nat.lt_mul_sub_one_div_div_add_two_mul}
  \leanok
  For $0 < p$ and $u < 2^{64}$ the Barrett reduction gives $u \bmod p$. For
  $p < \mathrm{maxPrimeEven}$, coefficients below $p$ and $a < p$, the accumulator of Horner's rule
  stays below $p^{\mathrm{hornerSteps} + 1} < 2^{64}$, so that the schedule gives $f(a) \bmod p$,
  in both of its branches.
\end{lemma}

\begin{proof}
  \leanok
  Write $2^{64} - 1 = m p + t$ with $t < p$. Then $u/p - u m/2^{64} = u (1 + t)/(p\, 2^{64})$ lies
  in $[0, 1)$, so the quotient $\lfloor u m / 2^{64} \rfloor$ is $\lfloor u / p \rfloor$ or one
  less, and one subtraction of $p$ ends in $[0, p)$. After a reduction the accumulator is below $p$,
  and each step $s \mapsto s a + c$ with $a, c < p$ takes a bound $p^j$ to $p^{j + 1}$; a reduction
  follows at the latest after $\mathrm{hornerSteps}$ steps, so the accumulator stays below
  $p^{64 / \mathrm{maxBitsInPrime}}$, which is less than
  $2^{\mathrm{maxBitsInPrime} \cdot (64 / \mathrm{maxBitsInPrime})} \le 2^{64}$, whatever the
  degree. This is the range obligation of the accumulator, which in the C code is an
  \verb|unsigned long|; the invariant of the loop holds the residue of the accumulator and this
  bound.
\end{proof}

\begin{lemma}[The table with the exit is the table]\label{lem:examinePrime_eq}
  \uses{lem:barrett, def:examinePrimeExit, def:examinePrime}
  \lean{Ratpoints.Sieve.ExaminePrimeInvariant, Ratpoints.Sieve.examinePrime?_spec, Ratpoints.Sieve.run_examinePrime?, Ratpoints.Sieve.eq_examinePrime_of_run_examinePrime?_eq_some, Ratpoints.Sieve.eq_examinePrime_of_lt_numPrimes_of_run_examinePrime?_eq_some, Ratpoints.Sieve.eq_map_range_of_run_examinePrime?_eq_some, Ratpoints.Sieve.run_examinePrime?_eq_none_iff}
  \leanok
  The table of a prime with the exit (Definition~\ref{def:examinePrimeExit}) is a component of type
  \verb|Id| with the loop of \verb|examine_prime| over the residues: the entry $0$ from the squares
  modulo $p$ at $c_0$, the entries $1, \dots, p - 1$ by the schedule of
  Definition~\ref{def:barrett}, and the entry $p$, whether there is a point at infinity modulo $p$;
  it keeps $p + 1$ entries, as \verb|examine_prime| does, and hands back the coefficients modulo $p$
  with the table. The reference (Definition~\ref{def:examinePrime}) has the entry at infinity too:
  for even $d$ whether $c_d$ is a square modulo $p$, and true for odd $d$. For
  $0 < p < \mathrm{maxPrimeEven}$, in particular for the primes of the table at the positions below
  $\mathrm{numPrimes}$: the table with the exit is none exactly when no entry of the reference is
  true, and otherwise it is the reference, with the coefficients reduced modulo $p$.
\end{lemma}

\begin{proof}
  \leanok
  Lemma~\ref{lem:barrett} for each entry $1, \dots, p - 1$, in an invariant of the loop over the
  residues; the entry $0$ and the entry at infinity read the same squares in both.
\end{proof}

The lemmas on the rows and the exit (Lemmas~\ref{lem:admits_examinePrime}
and~\ref{lem:examinePrime_none}) are those of the table of a prime, the exit with the hypothesis
$p < \mathrm{maxPrimeEven}$, which every prime that the model examines satisfies, and with the lemma
that a point makes some entry of the reference true.

\section{The loop of the tables}\label{sec:examinePrimes}

\begin{definition}[The loop of the tables]\label{def:examinePrimes}
  \uses{def:examinePrimeExit, def:forbidden}
  \lean{Ratpoints.Sieve.examinePrimes}
  \leanok
  For the first $n$ primes of the table in turn: the table of the prime with the exit, and none when
  the curve has no points modulo it (\csource{find_points.c}{2605}{2614}); the table of every
  examined prime kept by its position, where \verb|sieving_info| keeps by its index the sieve entry
  of each prime that says something, with the codes $1$ and $2$
  (\csource{find_points.c}{2537}{2537}, \csource{find_points.c}{2545}{2545},
  \csource{find_points.c}{2601}{2602}, \csource{find_points.c}{2616}{2616},
  \csource{find_points.c}{2619}{2619}, \csource{find_points.c}{2622}{2622},
  \csource{find_points.c}{2707}{2707}); and, while the checked loop is on and fewer than
  $\mathrm{maxForbidden}$ forbidden divisors have been found, the forbidden divisor of the prime
  from the coefficients modulo $p$ and the entry at infinity that its table hands back
  (\csource{find_points.c}{2624}{2626}, \csource{find_points.c}{2677}{2677}). The unit parts, the
  forbidden valuations and the forbidden divisor of a prime
  (Definitions~\ref{def:forbiddenValuations} and~\ref{def:forbidden}) take these as arguments; the
  forbidden divisors of the further primes, the second loop of Definition~\ref{def:forbidden}, start
  from the lists of this loop at the position $n$, which stands for the final value of
  \verb|pn_lim|.
\end{definition}

\begin{lemma}[The loop of the tables]\label{lem:examinePrimes}
  \uses{def:examinePrimes, lem:examinePrime_eq, lem:examinePrime_none, def:forbidden, def:barrett}
  \lean{Ratpoints.Sieve.ExaminePrimesInvariant, Ratpoints.Sieve.ExaminePrimesLoopInvariant, Ratpoints.Sieve.examinePrimes_spec, Ratpoints.Sieve.size_eq_of_run_examinePrimes_eq_some, Ratpoints.Sieve.getElem_eq_examinePrime_of_run_examinePrimes_eq_some, Ratpoints.Sieve.isForbiddenEntry_of_run_examinePrimes_eq_some, Ratpoints.Sieve.isForbiddenVal_of_run_examinePrimes_eq_some, Ratpoints.Sieve.not_isPoint_of_run_examinePrimes_eq_none}
  \leanok
  For $n \le \mathrm{numPrimes}$: if the loop gives tables, there are $n$ of them, the table at the
  position $i$ is the reference table of the $i$-th prime of the table, and the forbidden divisors
  found are among those that the loop may record (the predicates of
  Definition~\ref{def:forbidden}), at most $\mathrm{maxForbidden}$ of them; if it gives none, some
  prime among the first $n$ has no table with the exit, and the curve has no points.
\end{lemma}

\begin{proof}
  \leanok
  An invariant of the loop over the primes seen so far, and one of the loop with its early return,
  which holds the result only at the end, as for the set-up of the Jacobi test. The primes of the
  table at positions below $\mathrm{numPrimes}$ are below $\mathrm{maxPrimeEven}$
  (Definition~\ref{def:barrett}), so Lemma~\ref{lem:examinePrime_eq} applies to each.
\end{proof}

\section{Prime powers}\label{sec:powers}

\begin{definition}[The inverse modulo $m$ by the extended Euclidean algorithm]\label{def:modinv}
  \lean{Ratpoints.Sieve.modinv}
  \leanok
  From $r_0 = m$, $r_1 = x \bmod m$ and the cofactors $s_0 = 0$, $s_1 = 1$, which keep
  $r_i \equiv s_i x \pmod m$: while $r_1 \ne 0$, with $q = \lfloor r_0 / r_1 \rfloor$, the pair
  $(r_0, r_1)$ becomes $(r_1, r_0 - q r_1)$ and $(s_0, s_1)$ becomes $(s_1, s_0 - q s_1)$; then
  $s_0$, plus $m$ if it is negative (\csource{find_points.c}{2241}{2252}).
\end{definition}

\begin{lemma}[The inverse modulo $m$]\label{lem:modinv}
  \uses{def:modinv}
  \lean{Ratpoints.Sieve.ModinvInvariant, Ratpoints.Sieve.modinv_spec, Ratpoints.Sieve.ModinvInvariant.natAbs_add_mul_le, Int.natAbs_sub_mul_of_mul_nonpos}
  \leanok
  For $1 < m$ and $x$ prime to $m$, the result is the inverse of $x$ modulo $m$ in $[0, m)$; the
  cofactors and the products
  $q s_1$ stay in $[-m, m]$, which is their range obligation.
\end{lemma}

\begin{proof}
  \leanok
  The invariant: the two congruences, $\gcd(r_0, r_1) = \gcd(m, x)$, $r_0 > 0$, cofactors of
  opposite signs with $|s_0| r_1 + |s_1| r_0 = m$, and $|s_0| \le m$; the loop ends since $r_1$
  decreases. At the end $r_0 = \gcd(m, x) = 1$, so $s_0 x \equiv 1$.
\end{proof}

\begin{definition}[The tables of a prime power]\label{def:examinePower}
  \uses{def:modinv, def:useCLong, def:examinePrime}
  \lean{Ratpoints.Sieve.Power, Ratpoints.Sieve.examinePower, Ratpoints.Sieve.Power.admits}
  \leanok
  For a prime $p$ and $e \ge 2$, $q = p^e$, a component of type \verb|Id|: the coefficients reduced
  modulo $q$ on the path of the flag of the \verb|long|s, with $c_D = 0$ for odd $d$; the squares
  modulo $q$; for each residue $x$ whether $f(x)$ is a square modulo $q$, by Horner's rule reduced
  at every step; for each residue $t$ whether $p \mid t$ and $\tilde f(t)$ is a square modulo $q$;
  and, in a loop over the residues, the inverses of the units modulo $q$ by
  Definition~\ref{def:modinv} (\csource{find_points.c}{2267}{2303};
  \csource{rp-private.h}{523}{523}, \csource{rp-private.h}{531}{534}). The row of $b$ admits $a$
  when, for $r = b \bmod q$ and $x = a \bmod q$, $f(x r^{-1})$ is a square modulo $q$ if
  $p \nmid r$, and $p \nmid x$ and $\tilde f(r x^{-1})$ is a square modulo $q$ if $p \mid r$
  (\csource{init.c}{359}{386}, \csource{rp-private.h}{592}{592}).
\end{definition}

\begin{lemma}[The rows of a prime power let every point through]\label{lem:power_admits}
  \uses{def:examinePower, lem:modinv, lem:isSquare_div, lem:twoAdic_point, def:examinePrime,
    def:barrett, def:binaryForm}
  \lean{Ratpoints.Sieve.ExaminePowerInvariant, Ratpoints.Sieve.examinePower_spec, Ratpoints.Sieve.examinePower_eq_examinePower_false, Ratpoints.Sieve.getElem!_fsq_run_examinePower_of_isSquare, Ratpoints.Sieve.getElem!_gsq_run_examinePower_of_isSquare, Ratpoints.Sieve.mul_getElem!_inv_run_examinePower_mod, Ratpoints.Sieve.natCast_evalMod_reverse, Ratpoints.Sieve.admits_run_examinePower_of_sq_eq, Ratpoints.mul_add_lt_two_pow_64_of_lt_maxPrimeEven, Ratpoints.mul_self_lt_two_pow_63_of_lt_maxPrimeEven, Ratpoints.Sieve.getElem!_inv_run_examinePower_lt, Ratpoints.Sieve.getElem!_inv_run_examinePower_lt_two_pow_16}
  \leanok
  The two paths of the coefficients give the same tables; the tables mark what they claim: $f(x)$ a
  square modulo $q$, $\tilde f(t)$ a square modulo $q$ for $p \mid t$, and, by the contract of the
  loop of the inverses, the inverses of the units. For a prime $p$, $e \ge 1$, $\gcd(a, b) = 1$ and
  $y^2 = F(a, b)$, the row of $b$ in the tables of $p^e$ admits $a$. The values of
  \verb|examine_power| and of the row stay in their C types, the moduli being below
  $\mathrm{maxPrimeEven} \le 2^{10}$: the accumulator $v x + c$ of Horner's rule with $v, x, c < q$,
  an \verb|unsigned long|, and likewise the product of a residue and an inverse in the row; the
  square $x \cdot x$, a signed \verb|long|; and the inverses, below $q$, \verb|unsigned short|s.
\end{lemma}

\begin{proof}
  \uses{lem:modinv}
  \leanok
  For $p \nmid b$, $b$ is a unit modulo $q$ and $F(a, b) = b^D f(a b^{-1})$ there, with $b^D$ the
  square of a unit (Lemma~\ref{lem:isSquare_div}). For $p \mid b$, $p \nmid a$, so $a$ is a unit,
  and $F(a, b) = a^D \tilde f(b a^{-1})$ modulo $q$, by the form of that lemma with the roles of $a$
  and $b$ exchanged (Lemma~\ref{lem:twoAdic_point}); Horner's rule on the reversed coefficients
  computes $\tilde f$. Zero counts as a square, so a high power of $p$ dividing $F(a, b)$ needs no
  case of its own.
\end{proof}

The coefficients reduced by \verb|mod| need moduli below $2^{59}$ (Definition~\ref{def:useCLong}),
which $\mathrm{maxPrimeEven} \le 2^{32}$ gives.

\section{The sieve entry of a modulus}\label{sec:sieveEntry}

\begin{definition}[The factors of a modulus]\label{def:modulusFactors}
  \uses{def:primeTable}
  \lean{Ratpoints.modulusFactors}
  \leanok
  The factors of $m$ over the first $\mathrm{numPrimes}$ primes of the table, as \verb|add_moduli|
  finds them: for the primes $p$ in turn while the cofactor is above $1$, a cofactor below $p^2$,
  which has no odd prime factor below $p$ and so, for odd $m$, is a prime, is looked up among the
  primes from $p$ on and taken with the exponent $1$ when it is one of them; else $p$ is divided off
  as often as it divides. The result is the list of the positions and exponents, in the order of
  the table, and the cofactor that is left (\csource{find_points.c}{2362}{2362},
  \csource{find_points.c}{2365}{2365}, \csource{find_points.c}{2369}{2386},
  \csource{find_points.c}{2396}{2400}). The factorization is complete exactly when the cofactor is
  $1$; an even modulus never factors completely, as the table holds no $2$, and \verb|add_moduli|
  offers odd moduli only.
\end{definition}

\begin{lemma}[The factors of a modulus]\label{lem:modulusFactors}
  \uses{def:modulusFactors}
  \lean{Ratpoints.ModulusFactorsInvariant, Ratpoints.ModulusFactorsPowInvariant, Ratpoints.modulusFactors_spec, Ratpoints.eq_mul_prod_run_modulusFactors, Ratpoints.lt_numPrimes_of_mem_run_modulusFactors, Ratpoints.one_le_of_mem_run_modulusFactors, Ratpoints.pow_dvd_of_mem_run_modulusFactors, Ratpoints.getElem!_primeTable_le_of_le, Ratpoints.getElem!_primeTable_inj, Ratpoints.getElem!_primeTable_ne_snd_run_modulusFactors, Ratpoints.snd_run_modulusFactors_lt_mul_self_of_dvd, Ratpoints.snd_run_modulusFactors_eq_one, Ratpoints.fst_run_modulusFactors_getElem!_primeTable}
  \leanok
  $m$ is the cofactor times the product of the powers of the primes at the positions; the positions
  increase and are below $\mathrm{numPrimes}$, the exponents are at least $1$, the cofactor is none
  of the first $\mathrm{numPrimes}$ primes of the table, and each of them that divides it is above
  its square root. So each factor divides $m$, and if every divisor $q > 1$ of $m$ is divisible by
  one of the first $\mathrm{numPrimes}$ primes of the table, as for those primes themselves, the
  cofactor is $1$ and the product of the factors is $m$; one of these primes is its only factor,
  with the exponent $1$, since the table increases, so that different positions hold different
  primes.
\end{lemma}

\begin{proof}
  \leanok
  An invariant of the loop over the primes and one of the inner loop that divides a prime off,
  which ends since the cofactor decreases. No prime passed divides the cofactor; at a stop below
  $p^2$ the lookup shows that the cofactor is none of the primes from $p$ on, so a cofactor above
  $1$ would have a divisor (itself, or $2$ when it is even) that no prime of the table divides.
\end{proof}

\begin{definition}[The sieve entries]\label{def:SieveEntry}
  \uses{def:admits, def:examinePower, def:binaryForm}
  \lean{Ratpoints.Sieve.Factor, Ratpoints.Sieve.Factor.admits, Ratpoints.Sieve.SieveEntry, Ratpoints.Sieve.SieveEntry.admits, Ratpoints.Sieve.survives, Ratpoints.Sieve.SieveEntry.IsSound}
  \leanok
  A factor of a modulus is a prime with its table or a prime power with its tables; the sieve entry
  of a modulus $m$ is $m$ with its factors, one for a prime or a prime power, two or more for a
  product, the reciprocal of $m$ (Definition~\ref{def:mulMod}), and the fields that the packing of
  the numerators reads (Chapter~\ref{chap:packing}): the bias, the least multiple of $m$ from
  $\mathrm{rowBias}$ on, the inverse $\mathrm{rbainv}$ of $\mathrm{rbaLength}$ modulo $m$, and the
  inverses $\mathrm{dinv}$ of the strides $2^k$ modulo $m$ up to the largest stride in use, which
  the rows of offsets fill (Definition~\ref{def:classOffsets}; \csource{rp-private.h}{497}{498},
  \csource{rp-private.h}{520}{520}, \csource{rp-private.h}{542}{555}). Its row of $b$ admits $a$
  when the row of every factor admits $a$, each factor reading $b$ and $a$ modulo its own modulus,
  which divides $m$ (\csource{init.c}{399}{414}, \csource{init.c}{419}{429},
  \csource{rp-private.h}{593}{593}); a numerator \emph{survives} a list of sieve entries for $b$
  when the row of $b$ of every entry admits it, the test that Definition~\ref{def:admits} gave for
  the tables of primes. A sieve entry is \emph{sound} when its modulus is positive and its row of
  $b$ admits $a$ for every $\gcd(a, b) = 1$ and $y^2 = F(a, b)$.
\end{definition}

\begin{definition}[The sieve entry of a modulus]\label{def:makeModulus}
  \uses{def:SieveEntry, def:modulusFactors, def:examinePrimeExit, def:examinePower, def:mulMod,
    def:packingConstants, def:rbaPack}
  \lean{Ratpoints.makeModulus?}
  \leanok
  A component of type \verb|Id|: by the factors of $m$, a factor with the exponent $1$ takes the
  table of its prime from the loop of the tables when its position is below the number examined, and
  is examined now, with the exit, otherwise; a factor with the exponent $e \ge 2$ gets the tables of
  its prime power; the entry holds the reciprocal of $m$, the bias
  $m \lceil \mathrm{rowBias} / m \rceil$ and the inverse $\mathrm{rbainv}$ of $\mathrm{rbaLength}$
  modulo $m$, with the formulas by which \verb|examine_prime| sets them for a prime, reading
  $\mathrm{rbainv}$ from the table of inverses at the remainder of $\mathrm{rbaLength}$ by the
  reciprocal, \verb|make_power| for a prime power and \verb|make_modulus| for a product, computing
  it by Definition~\ref{def:modinv}, and leaves the inverses $\mathrm{dinv}$ empty for the rows of
  offsets (Definition~\ref{def:classOffsets}) to fill (\csource{find_points.c}{2414}{2429},
  \csource{find_points.c}{2434}{2475}, \csource{find_points.c}{2546}{2546}; the tables of a prime
  power that \verb|add_moduli| examines when a candidate first needs it and keeps by its index,
  \csource{find_points.c}{2350}{2354}, \csource{find_points.c}{2358}{2359},
  \csource{find_points.c}{2387}{2395}, with their count, \csource{find_points.c}{2548}{2548}; the
  primes examined for the third stage, \csource{find_points.c}{2976}{2978},
  \csource{find_points.c}{2994}{2998}).
\end{definition}

\begin{lemma}[The rows of a modulus let every point through]\label{lem:sieveEntry_admits}
  \uses{def:makeModulus, lem:modulusFactors, lem:admits_examinePrime, lem:power_admits,
    lem:examinePrime_eq, lem:examinePrime_none, def:mulMod, lem:admits_mul, lem:modinv, lem:mulMod,
    def:packingConstants}
  \lean{Ratpoints.Sieve.Factor.modulus, Ratpoints.Sieve.Factor.IsSound, Ratpoints.MakeModulusInvariant, Ratpoints.IsModulusEntry, Ratpoints.mod_two_eq_one_of_forall_dvd, Ratpoints.coprime_rbaLength_of_mod_two_eq_one, Ratpoints.factors_eq_of_run_makeModulus?_eq_some, Ratpoints.makeModulus?_spec}
  \leanok
  For tables of the loop that are those of the reference and a modulus
  $1 < m < \mathrm{maxPrimeEven}$ each of whose divisors $q > 1$ is divisible by one of the first
  $\mathrm{numPrimes}$ primes of the table, as the well-formedness of the parameters asks of their
  moduli: if the entry of $m$ is made, it is an entry of $m$ whose row of $b$ admits $a$ for every
  $\gcd(a, b) = 1$ with $y^2 = F(a, b)$, whose factors have moduli that divide $m$ and tables of
  inverses that invert their units (Lemma~\ref{lem:admits_mul}), and which is \emph{packable}
  (Lemma~\ref{lem:classOffsets}): $m$ is odd, the entry holds the reciprocal of $m$,
  $\mathrm{rbainv} < m$ with $\mathrm{rbainv} \cdot \mathrm{rbaLength} \equiv 1 \pmod m$, and the
  bias is a multiple of $m$ in $[\mathrm{rowBias}, \mathrm{rowBias} + m)$; and for one of the first
  $\mathrm{numPrimes}$ primes of the table its only factor is that prime with the table of the
  reference; if it is not made, the curve has no points.
\end{lemma}

\begin{proof}
  \uses{lem:stage3, lem:classOffsets}
  \leanok
  Each factor's modulus divides $m$, so $b \bmod m$ and $a \bmod m$ reduce to the residues of the
  factor (Lemma~\ref{lem:admits_mul}); each factor admits the point by
  Lemma~\ref{lem:admits_examinePrime} or Lemma~\ref{lem:power_admits}, and its table of inverses
  inverts its units, by the inverses modulo a prime and by Lemma~\ref{lem:power_admits} for a prime
  power. A prime factor beyond the loop has its position below $\mathrm{numPrimes}$
  (Lemma~\ref{lem:modulusFactors}), so it is below $\mathrm{maxPrimeEven}$ and the table that its
  examination gives is the reference (Lemma~\ref{lem:examinePrime_eq}). No sieve entry is made only
  when a prime examined for it has no table with the exit, and then the curve has no points
  (Lemma~\ref{lem:examinePrime_none}). No Chinese remainder theorem is needed: the entry only
  intersects the conditions of its factors. A prime among the first $\mathrm{numPrimes}$ of the
  table is its own only factor (Lemma~\ref{lem:modulusFactors}), whose table the loop has kept or
  the examination gives as the reference (Lemma~\ref{lem:examinePrime_eq}). The table holds no $2$,
  so $m$ is odd and prime to $\mathrm{rbaLength} = 2^{\mathrm{rbaShift}}$; the factorization of $m$
  is complete, so the moduli of its factors multiply to $m$ (Lemma~\ref{lem:modulusFactors}): the
  single factor of a prime is the table of $m$, whose inverse at $\mathrm{rbaLength} \bmod m$, which
  the reduction with the reciprocal computes since $\mathrm{rbaLength}\, m < 2^{64}$
  (Lemma~\ref{lem:mulMod}), is the inverse of $\mathrm{rbaLength}$, below $m$ as every inverse of
  the table (Lemma~\ref{lem:stage3}), and for a prime power or a product Definition~\ref{def:modinv}
  computes it (Lemma~\ref{lem:modinv}). The bias $m \lceil \mathrm{rowBias} / m \rceil$ lies in
  $[\mathrm{rowBias}, \mathrm{rowBias} + m)$ by the division with remainder.
\end{proof}

The rows of a sieve entry are made from its admission test by the plain builder of
Definition~\ref{def:row}, for every kind; the bits of a row depend on the entry only through the
test, which reduces $b$ and $a$ modulo $m$ first, so the words of Chapter~\ref{chap:words} and their
contracts hold for the sieve entries of every kind with the proofs that they had for the tables of
primes. The builders of ratpoints, the functions of \verb|init.c| for each prime,
\verb|lay_out_pattern| for the prime powers and the blocks of \verb|_ratpoints_sieve_init_product|,
come with stage~F4.

\section{The parameters and the search}\label{sec:tablesSearch}

The parameters (Definition~\ref{def:Params}) have the number of primes that the loop of the tables
examines, from which the forbidden divisors of the further primes start; their moduli may be primes
of the table, powers of them and products of such powers. They are well-formed when every modulus
$m$ has $1 < m < \mathrm{maxPrimeEven}$ and factors over the first $\mathrm{numPrimes}$ primes of
the table, stated as: every divisor $q > 1$ of $m$ is divisible by one of them; when the moduli of
the third stage, from the position $s_2$ on, are primes of the table; when at most
$\mathrm{numPrimes}$ primes are examined; and when $s_1 \le s_2 \le$ the number of moduli. The
divisor form says what the factorization of Definition~\ref{def:modulusFactors} finds
(Lemma~\ref{lem:modulusFactors}) without its loop, so that the statements that comparator certifies
can be read without it. For the values of the constants of the build these are the odd moduli from
$3$ to $255$, the first $\mathrm{numPrimes}$ primes being the odd primes below $256$. Pairwise
coprimality, the bound of three factors, the bound $\mathrm{RATPOINTS\_COMPOSITE\_MAX}$, the cost of
the table, the room for the tables of the prime powers, what a modulus says about the curve and the
number of candidates are rules of the choice of the moduli, which soundness does not need. The
parameters of the command line examine the first $q$ primes and take primes as moduli
(Definition~\ref{def:sieveParams}).

\begin{definition}[What \texttt{sieving\_info} computes]\label{def:sievingInfo}
  \uses{def:examinePrimes, def:forbidden, def:makeModulus, def:Params}
  \lean{Ratpoints.sievingInfo}
  \leanok
  The loop of the tables over the examined primes (Definition~\ref{def:examinePrimes}); for the
  checked loop on, the forbidden divisors of the further primes; the sieve entries of the moduli,
  in a loop over them in their order (Definition~\ref{def:makeModulus}); and the array of the
  reciprocals of the entries, in their order (\csource{find_points.c}{3064}{3070}); none when the
  curve has no points modulo a prime examined on the way (\csource{find_points.c}{2522}{2532},
  \csource{find_points.c}{2892}{2900}, \csource{find_points.c}{3102}{3103}).
\end{definition}

\begin{lemma}[What \texttt{sieving\_info} computes is sound]\label{lem:sievingInfo}
  \uses{def:sievingInfo, lem:examinePrimes, lem:sieveEntry_admits, def:forbidden, def:mulMod,
    def:Params, def:stage3, lem:stage3, lem:classOffsets}
  \lean{Ratpoints.SievingInfoInvariant, Ratpoints.sievingInfo_spec}
  \leanok
  For well-formed parameters and either value of the switch of the checked loop: if
  \verb|sieving_info| gives a result, the sieve entries are those of the moduli, in their order,
  every sieve entry is sound, packable, so that it holds the reciprocal of its modulus, and has
  factors whose tables of inverses invert their units (Lemma~\ref{lem:sieveEntry_admits}), those of
  the third stage, from the position $s_2$ on, are prime entries (Definition~\ref{def:stage3}), the
  reciprocals returned are those of the entries in their order, and the forbidden divisors are among
  those that the two loops may record; if it gives none, the curve has no points.
\end{lemma}

\begin{proof}
  \leanok
  Lemma~\ref{lem:examinePrimes} gives tables equal to the reference, which
  Lemma~\ref{lem:sieveEntry_admits} needs, in an invariant of the loop over the moduli, and the
  conditions on a modulus that it takes come from the well-formedness; the forbidden divisors of the
  second loop are those of the contract of Definition~\ref{def:forbidden}. The exit is that of the
  loop of the tables or of a modulus. The reciprocal of each entry is that of its modulus by
  Lemma~\ref{lem:sieveEntry_admits}. A modulus of the third stage is one of the first
  $\mathrm{numPrimes}$ primes of the table (Definition~\ref{def:Params}), whose entry has that prime
  with the table of the reference as its only factor (Lemma~\ref{lem:sieveEntry_admits}) and its
  reciprocal, a prime entry since the prime is below $2^{32}$ (Lemma~\ref{lem:stage3}).
\end{proof}

The search (Definition~\ref{def:searchPoints}) takes its sieve entries and its forbidden divisors
from \verb|sieving_info|, and otherwise runs as in Chapter~\ref{chap:denominators}; the soundness of
the run over a shape takes sound sieve entries. In this sub-stage its theorem
(Theorem~\ref{thm:findPoints_order}), the main theorems and the refusal theorems keep their text:
the output is the list of the points at infinity followed by the points of the brute-force reference
over the denominators of the shape in its order, each mapped back after a reversal, for every
well-formed choice of the parameters, also when an exit is taken.

\section{The command-line program and the differential tests}\label{sec:tablesCli}

The program is that of stage~E: its parameters examine the first $q$ primes and take primes as
moduli. F1 makes no further invocation of the test scripts comparable, since the options that it
concerns have been read since stage~C1 and the output does not change; the differential tests stay
as they are, every point list compared in its order. Composite moduli, which ratpoints chooses from
the curve, reach the model through the parameters of the theorems; \verb|#guard| examples run the
search with composite moduli in the stages, with the number of primes examined below and above the
primes of the moduli, and compare the output with the brute-force reference, and with the reference
of the reversed polynomial mapped back when the search reverses.

\section{Deviations}\label{sec:tablesDeviations}

Each component of this chapter differs from the C code in a stated way, recorded in its
\verb|Deviation:| paragraph:
\begin{itemize}
  \item the arithmetic of \verb|unsigned long|s in the Barrett reduction and in Horner's rule is
    that of natural numbers, with the bound $2^{64}$ of the accumulator a lemma
    (Lemma~\ref{lem:barrett}); the branch of the schedule is chosen for each residue, where
    \verb|examine_prime| chooses it once; natural numbers for the remainders of \verb|modinv|, whose
    correction of a negative $x \bmod m$ then has no counterpart, and for the \verb|long|s of
    \verb|examine_power| and of the rows of prime powers, integers for the cofactors of
    \verb|modinv|, with their bounds lemmas (Lemmas~\ref{lem:modinv} and~\ref{lem:power_admits});
  \item \verb|examine_prime| counts the residues with a square value; the model asks whether an
    entry of the table is set, which is the same test. Its code $0$, every residue a square and no
    point at infinity, makes no sieve entry in ratpoints and an entry in the model, which the
    parameters may use as a modulus: its rows let every numerator through for a unit $b$. The fields
    of the sieve entry of a prime, the reciprocal, the bias and the inverse of $\mathrm{rbaLength}$
    among them, are those of a sieve entry with the prime as its single factor, set where the model
    makes the entry of a modulus;
  \item the coefficients and the tables are fresh arrays, where the C code fills arrays of its
    caller and buffers; the tables of a prime power are computed for each modulus that needs them,
    where \verb|add_moduli| examines each prime power once and keeps it, and \verb|make_modulus|
    keeps the sieve entry of each prime power;
  \item the factors of a modulus are a list, for the codes of the factors and their number, without
    the bound of three factors or of the position $64$, which serve the choice of the moduli, and
    over the first $\mathrm{numPrimes}$ primes, where \verb|add_moduli| factors over the examined
    ones; a prime factor beyond the examined primes is examined when its modulus is made, where
    ratpoints examines further primes in the order of the table, in the loop of the tables while its
    first two stages want more primes and for its third stage, both only when the number of primes
    was left to the library; a factor whose prime says nothing takes its table from the loop, where
    \verb|add_moduli| offers no composite with such a factor, and a prime-power factor whose tables
    say nothing is used with them, where \verb|add_moduli| refuses such a modulus;
  \item the rows of a prime power and of a product are made by the plain builder from the admission
    test, where \verb|_ratpoints_sieve_init_power| steps through the residues and lays the pattern
    out with \verb|lay_out_pattern|, and \verb|_ratpoints_sieve_init_product| lays down the rows of
    its factors in blocks, read with an acquire load and published with a release store
    (stage~F4);
  \item the density, the ranking, the choice of the moduli and their stages are not modelled
    (stage~I): the moduli, their stages and the number of primes examined are those of the
    parameters, and the parameters may hold moduli that ratpoints never offers (composites above
    $\mathrm{RATPOINTS\_COMPOSITE\_MAX}$, with a factor beyond the examined primes, or with a factor
    that says nothing), all sound; the parameters of the command line without \verb|-p| take
    primes after the first $q$ when the counts of their stages ask for more, primes that ratpoints
    then may examine too, and the model makes their entries with the others;
  \item the reciprocal, the bias and the inverse of $\mathrm{rbaLength}$ of a sieve entry are set
    with the entry for every kind of modulus, with the formulas of the C code
    (Definition~\ref{def:makeModulus}), the inverse of a prime read from its table; the inverses
    \verb|dinv| an array that the rows of offsets fill (Definition~\ref{def:classOffsets}).
\end{itemize}
The lines of this part of the C code that are neither modelled nor left to a later stage are
excluded in \verb|correspondence/ratpoints.toml|: the check \verb|RP_VERIFY_MODULI| of the rows of
the moduli against the form, compiled only on request (\verb|init.c|); the output of
\verb|sieving_info| under \verb|RP_PRIME_STATS|, \verb|RATPOINTS_VERBOSE| and \verb|DEBUG| between
statements that are not modelled, where the output inside a modelled block goes with the block,
whose deviation says so; and the enlargement of the buffer of the tables for every prime looked at,
those of the third stage among them, and for the composite moduli of the first two stages with their
prime-power factors, which fails only for lack of memory. The reciprocals of the moduli are fields
of their sieve entries, which \verb|sieving_info| collects at its end
(\csource{find_points.c}{3064}{3070}; Definition~\ref{def:sievingInfo}); the third stage reads them
(Definition~\ref{def:stage3}), and so do the first two stages, stepping them with the sieve specs
(Definitions~\ref{def:andRows} and~\ref{def:secondStage}); left to stage~I are the density of a
prime and of a power, \verb|add_moduli| beyond the factorization, the ranking and the choice of the
moduli with the rule for looking at further primes.

\section{Files and names}\label{sec:tablesFiles}

The constants $\mathrm{maxPrimeEven}$ and $\mathrm{hornerSteps}$ are in
\verb|Ratpoints/Model/Constants.lean|, each citing its macro, computed from
$\mathrm{maxBitsInPrime}$ so that the build with other values of the constants does not vary them
itself, with the two facts decided there, $1 \le \mathrm{hornerSteps}$ and
$\mathrm{maxBitsInPrime} \le 10$; the bounds $\mathrm{maxPrimeEven} \le 2^{10}$ and $\le 2^{32}$ are
in \verb|Ratpoints/Proof/Reductions.lean|, without Mathlib, and the other bounds derived from them
in \verb|Ratpoints/Proof/Params.lean|. The tables of a prime with the Barrett reduction and Horner's
rule on its schedule, \verb|modinv|, the tables of a prime power, the factors and the sieve entries
are in \verb|Ratpoints/Model/Tables.lean|, in the namespace \verb|Ratpoints.Sieve|, with the squares
and the inverses, the helper \verb|mod| and the rows of a denominator of the earlier stages; the
factors of a modulus are in \verb|Ratpoints/Model/Params.lean|; the loop of the tables and the
second loop of the forbidden divisors in \verb|Ratpoints/Model/Denominators.lean|; and the sieve
entry of a modulus and \verb|sieving_info| in \verb|Ratpoints/Model/FindPoints.lean|. The contracts
are next to the lemmas they serve: the Barrett reduction, Horner's rule and the table with the exit
in \verb|Ratpoints/Proof/Tables.lean|; \verb|modinv|, the prime powers with the bounds of their
inverses, and the sieve entry of a modulus in \verb|Ratpoints/Proof/Moduli.lean|, with what the
packing reads of it in \verb|Ratpoints/Proof/Packing.lean|, without Mathlib; the facts on the table
of primes, the other range lemmas of \verb|examine_power|, the factors of a modulus and the
well-formedness in \verb|Ratpoints/Proof/Params.lean|; the loop of the tables and the forbidden
divisors in \verb|Ratpoints/Proof/Forbidden.lean|; the sieve entries in the contracts of
\verb|Ratpoints/Proof/Sieve.lean|, \verb|Words.lean|, \verb|Stages.lean| (the prime entries of the
third stage), \verb|Sift.lean| and \verb|Denominators.lean|; and the soundness of the sieve entries,
\verb|sieving_info| and the search in \verb|Ratpoints/Proof/FindPoints.lean|. The lemmas on words in
core's namespaces are in \verb|Ratpoints/Proof/Bits.lean|, those on arrays and integers in
\verb|Ratpoints/Proof/Containers.lean|, and the lemma on sorted lists in the mirror
\verb|Ratpoints/Mathlib/Data/List/Sort.lean|, as a lemma on \verb|List.SortedLE|.

\section{The evidence}\label{sec:tables-evidence}

Besides the proofs:
\begin{itemize}
  \item a reading of the C code of the numerator side at the pin, each part with a second reader who
    checked the reading: every part of the numerator side is a data refinement of the model or a
    sound filter, so that the points and their order do not change; the Barrett reduction is exact
    for every $u < 2^{64}$, and the schedule of Horner's rule keeps its accumulator below $2^{64}$
    for the primes below $2^{\mathrm{maxBitsInPrime}}$, at every degree; and two defects of the C
    code without effect on the points, now recorded with ratpoints: \verb|-N| without \verb|-n| is
    not a hard limit, and the check of \verb|RATPOINTS_CHUNK| above $16$ never fires;
  \item before the sub-stage, a sketch of the model in Lean, which agrees with \verb|ratpoints -q|
    on the 3219 runs of the differential tests and the 13 at heights up to 100000, every point list
    in its order; whose tables, coefficients modulo $p$, tables of prime powers and rows are those
    of the C code (33072 tables of primes, 3168 of prime powers, 5.2 million rows for every odd
    modulus from $3$ to $255$, and at $\mathrm{maxBitsInPrime} = 10$ against a build with that
    value); whose forbidden divisors are those of an instrumented build on 785 runs; and whose
    search with random composite moduli gives the output of the default parameters on 420 runs; a
    prototype of the proofs, which showed that soundness needs only each factor of a modulus to
    divide it, without the Chinese remainder theorem; and a review of the C ranges of the sketch (13
    findings, all confirmed) and of the plan of this chapter by four readers (38 findings, 30
    confirmed and 7 in part), each finding checked by a second reader;
  \item after the sub-stage, two independent reviews, each finding checked by a second reader: of
    the C ranges, the deviations and the exclusions against the C code (10 findings, all confirmed:
    the buffer of the prime powers to the sieve entry of a modulus, a dead store and the frames of
    the generators of the tables to the exclusions, and the deviations name the lemmas of their
    range obligations and what they leave out), and of the code against the conventions of the
    project (54 findings, 47 confirmed), applied in twelve rounds, among them the lemma on sorted
    lists in the currency of Mathlib, the shared facts as lemmas of their own, conjunctions split
    and unused lemmas removed, and the tables and the sieve entries in a module of their own;
  \item the \verb|#guard| examples of the model: the Barrett reduction at the ends of its range,
    Horner's rule on its schedule against Horner's rule reduced at every step, the tables and the
    coefficients modulo $p$ of \verb|examine_prime| and the tables of \verb|examine_power| read from
    a build of ratpoints, the inverses, the factors and the sieve entries of moduli of each kind,
    the exits, and the search with composite moduli against the brute-force reference;
  \item the differential tests of Section~\ref{sec:tablesCli}: 3219 comparisons in the default
    suites and 13 in \verb|test4high|, none failed, every point list compared in its order;
  \item Comparator accepted the solution of the challenge, with the parameters of this sub-stage, on
    2026-10-13;
  \item the project builds, with all its proofs, for other values of the constants of the build,
    $\mathrm{maxBitsInPrime} = 5$ and $10$ among them; so no proof depends on their values;
  \item every proof of the sub-stage takes at most 20 lines.
\end{itemize}
